\section{결론 및 향후 연구}\label{sec:conclusion}

본 논문은 범용 MCU 위 SNN 추론의 병목을 절제 실측으로 재진단하고,
그 결과에 따라 커널을 재설계하여 STM32F746에서 이미지당
$2{,}328.5$ms를 $380.6$ms로, \textbf{6.12배} 가속하였다. 재진단은
통념과 달리 명령 페치 최적화(ART, ITCM 코드 배치)가 이 워크로드에서
측정 노이즈 이내의 무효과임을 보였고, 진짜 병목이 ``0을 곱하기 위한
데이터 로드''임을 특정했다. 제안한 무곱셈 희소 커널은 이진 활성의 합
항등식~(\ref{eq:collapse})을 Cortex-M7의 메모리 계층에 맞게 --- im2col
제거, 전치 가중치 $[k][C_o]$, 8채널 레지스터 블로킹으로 --- 구현하여
발화율 4.1\%인 두 번째 합성곱을 동일 배치의 dense 대비 13.0배
가속했으며, requantize$\cdot$LIF
에필로그를 dense 경로와 공유함으로써 구성상 비트 동일성을 유지했다.
여섯 빌드 변형의 출력 체크섬 일치가 이를 실리콘에서 검증한다. 제안
커널은 TFLM 커스텀 연산자 프레임과 CMSIS-NN 지원 함수 위에 그대로
얹히며, 런타임 오버헤드는 2.1\%에 머문다.

이 결과의 일반적 시사는 두 가지다. 첫째, SNN의 저전력$\cdot$저지연
서사가 약속하는 이벤트 기반 이득은 뉴로모픽 하드웨어 없이도, 상용
MCU의 정수 커널 수준에서 대수적 항등식만으로 회수될 수 있다 ---
단, 그 이득의 상한은 발화율의 역수로 정해진다. 둘째, 마이크로벤치마크가
지목하는 병목(페치 경로)과 실제 워크로드의 병목(데이터 로드)은
다를 수 있으며, 절제 가능한 빌드 매트릭스와 비트 회귀 검증을 갖춘
사이클 정밀 하네스가 그 간극을 드러내는 실질적 도구였다. 다섯 번의
커널 반복 중 한 번은 실측 퇴행이었고(가중치 전치가 DTCM 상주
레이어에서 역효과), 이를 잡아낸 것도 같은 하네스다.

향후 과제는 다음과 같다. (i) 현재 최대 항목이 된 CPL1(밀도 47.7\%,
전체의 58\%)에 대해, 인코더가 스파이크 좌표 목록을 직접 방출하는
이벤트 기반 재구성으로 수용 영역 재스캔(입력의 스텝당 약 9회
재독)을 제거하는 것 --- 정수 합의 순서 무관성 덕에 비트 동일성은
유지 가능하다. (ii) FC 가중치의 SRAM 사본과 216MHz 승격 등 잔여
최적화. (iii) 발화율이 다른 모델$\cdot$코딩 방식(TTFS 등)에서 제안
커널의 이득 곡선을 일반화하는 것. (iv) 정수 구현의 전체 테스트셋
정확도를 보드에서 직접 실측하는 것은 선행 연구와 공유하는 한계로
남아 있다.
